\title{Parameter-Efficient Orthogonal Finetuning via Factorization}

\begin{abstract}
Large foundation models have become widespread, yet training them anew is prohibitively resource-intensive. Consequently, strategies for cost-effective adaptation of these robust models to specialized tasks have gained substantial relevance. Here, we focus on a systematic finetuning approach—Orthogonal Finetuning (OFT)—for adapting foundation models to downstream applications. While OFT provides strong generalization, it remains parameter-heavy owing to the substantial size of orthogonal matrices. To mitigate this, we analyze OFT from the lens of information transmission and delineate several critical criteria that contribute to improved parameter efficiency. Motivated by the efficiency of the Cooley-Tukey fast Fourier transform algorithm in information propagation, we introduce an optimized orthogonal parameterization utilizing butterfly structures. Integrating this parameterization into OFT yields a novel and more parameter-efficient adaptation strategy, termed Orthogonal Butterfly (BOFT). BOFT generalizes OFT by incorporating it as a specific instance within its broader framework. We support our proposal with an extensive set of experiments, demonstrating the adaptation of large vision transformers, extensive language models, and text-to-image diffusion architectures across a range of tasks in both vision and language domains.
\end{abstract}

\section{Introduction}

State-of-the-art models like ChatGPT~\cite{brown2020language,chen2021evaluating} and Stable Diffusion~\cite{rombach2022high} highlight the exceptional capacity for generalization characteristic of large foundation models. The substantial gains in their performance have been accompanied by an explosive growth in parameter counts (\emph{e.g.}, GPT-3 incorporates approximately 175 billion parameters). This tremendous scale renders it increasingly impractical for most researchers to train such models from the ground up. Thus, the field's ongoing progress depends on our capability to tailor pre-trained foundation models to new tasks efficiently, without full retraining. There are three main classes of methods for efficient downstream adaptation: \emph{model finetuning}~\cite{hulora2022,zaken2022bitfit,guo2020parameter,raffel2020exploring,qiu2023controlling,zhang2022adaptive,chavan2023one}, wherein the original architecture is preserved and only a portion of the parameters is updated; \emph{adapter tuning}~\cite{rebuffi2017learning,houlsby2019parameter,pfeiffer2020adapterfusion,lin2020exploring,he2021towards}, which adds small trainable modules to the base model; and \emph{prompt tuning}~\cite{li2021prefix,lester2021power}, where trainable prefix tokens augment the model input. Of these techniques, model finetuning stands out as a straightforward yet potent option, offering the key advantage of maintaining original inference speed.

The core idea underpinning model finetuning is to maintain a resemblance between the pretrained and finetuned models, according to specified metrics, thereby retaining the knowledge acquired during pretraining. Common finetuning strategies, such as using a low learning rate, effectively serve as heuristics to minimize the divergence between these two models. Considering the structured arrangement of weight matrices, a more systematic strategy aims to preserve the intrinsic relational properties of these matrices—specifically, the angles formed between pairs of neurons. Building on this observation, the Orthogonal Finetuning~(OFT) framework~\cite{qiu2023controlling} has been introduced, in which neurons within a layer are all subjected to the same orthogonal transformation, guaranteeing the conservation of their pairwise angular relationships throughout finetuning. While OFT has exhibited strong results in terms of generalization and convergence, particularly for text-to-image diffusion model finetuning~\cite{qiu2023controlling}, a significant challenge arises from the extensive number of parameters required for high-dimensional orthogonal matrices. To counteract this, OFT incorporates a block-diagonal configuration to lessen the parameter count. Nevertheless, this approach also imposes certain limitations: the orthogonal matrices must adhere to a fixed block-wise sparsity, with each block processed independently. Although parameter-efficient, this block-diagonal structure can introduce unwanted inductive biases (for example, restricting the expressiveness of the transformations and preventing approximation of standard linear operations). Furthermore, determining an optimal sparsity pattern for these blocks remains an unresolved issue.

A central challenge in this context is devising a dense orthogonal matrix in a manner that remains parameter-efficient. Typically, a $d$-dimensional dense orthogonal matrix necessitates $\mathcal{O}(d^2)$ parameters, which appears prohibitive for large $d$. We address this by constructing the dense orthogonal matrix through a composition of several factorized sparse matrices, thereby offering a new solution. This methodology capitalizes on the principle of reducing the parameter count by accepting an increased computational cost, essentially exchanging space for time. In this scheme, since the orthogonal matrix is realized as the product of sparse factors, repeated multiplications are required at every step of training. Reframing this matrix factorization, we conceptualize the generation of a dense orthogonal matrix as an exercise in information transmission. Under this paradigm, constructing a dense orthogonal matrix via sparse matrix products is analogous to propagating information across a grid-structured network. This analogy provides a systematic basis for designing various sparse matrix factorization techniques that limit the number of trainable parameters while maintaining the expressive capability necessary for producing dense matrices.

Our approach to enhancing parameter efficiency within this information transmission model draws heavily on the butterfly network configuration featured in the Cooley-Tukey fast Fourier transform algorithm~\cite{cooley1965algorithm}, renowned for facilitating rapid information exchange among nodes~\cite{klasing1994broadcasting}. In particular, the butterfly network of the Cooley-Tukey algorithm suggests a structured sparse factorization strategy—commonly referred to as butterfly factorization. Employing butterfly factorization allows us to construct a $d$-dimensional dense matrix by multiplying $\mathcal{O}(\log d)$ sparse matrices, where each contains only $\mathcal{O}(d)$ nonzero elements. By ensuring orthogonality for each sparse factor, we arrive at a highly parameter-efficient orthogonal representation that only requires $\mathcal{O}(d \log d)$ parameters, dramatically lowering parameter overhead compared to conventional full parameterizations. Leveraging this structure, we introduce a new, highly parameter-efficient finetuning scheme: Orthogonal Butterfly~(BOFT). Notably, BOFT encapsulates OFT as a specific instance and serves as a flexible framework for orthogonal finetuning. Both the block-diagonal and butterfly constructions share the property of sparsity, employing structured sparsity to reduce the total number of parameters subject to optimization. Drawing a parallel with LoRA~\cite{hulora2022}, our work emphasizes the distinction between parameter savings achieved through sparsity (as in our approach) and those accomplished with low-rank structure; both strategies fall under the umbrella of structured matrices~\cite{candes2011robust} designed to optimize parameter usage.

In contrast to the block-diagonal approach adopted by OFT, which attempts to balance model expressiveness with structural regularity, BOFT leverages the butterfly structure to establish a more continuous transition between full orthogonal matrices (enhancing expressiveness) and the identity matrix (promoting regularity). This approach facilitates the identification of a more constrained hypothesis space within the broader orthogonal group, tailored for specific downstream applications. The butterfly structure's extensive application in fast linear transformations—such as those found in the discrete Fourier and discrete cosine transforms—supports our assertion that enforcing this structure will impart a valuable inductive bias. Such a bias not only encourages better generalization but also mitigates the risk of overfitting. Our underlying reasoning is grounded in the butterfly structure's capacity to approximate a variety of canonical linear transforms, a feat unattainable by OFT's block-diagonal formulation. The principal contributions of our work are as follows:

\begin{itemize}[leftmargin=*,nosep]
    \item We address the challenge of parameter-efficient orthogonal finetuning by formulating an innovative information transmission framework. This formulation recasts the design of parameter-efficient dense orthogonal matrices as an equivalent information flow problem on grid-structured graphs.
    \item Drawing inspiration from the butterfly decomposition in the Cooley-Tukey algorithm, we introduce Orthogonal Butterfly—a parameter-efficient orthogonal finetuning technique—within this transmission paradigm.
    \item We present theoretical motivations elucidating how BOFT substantially curtails the count of trainable parameters without compromising expressive capacity or stability during optimization. The factorized representation in BOFT also inherently enables compelling forms of weight interpolation (refer to Figure~\ref{fig:boft_interpolation}).
    \item For the first time, we extend orthogonal finetuning methods beyond their established application in controllable text-to-image synthesis~\cite{qiu2023controlling}, highlighting their broader utility as universal model finetuning strategies. To substantiate this, we evaluate BOFT across multiple adaptation tasks spanning computer vision and natural language processing, where it surpasses existing state-of-the-art approaches, thus affirming its advancements in both parameter efficiency and generalization.
\end{itemize}

\textbf{Parameter-efficient finetuning (PEFT)}. As foundation models (\emph{e.g.}, \cite{brown2020language,radford2021learning,rombach2022high,kirillov2023segment}) scale in size and capacity, there has been growing focus on techniques that enable efficient adaptation for specific tasks with minimal changes to model parameters~\cite{treviso2023efficient,ding2023parameter,lialin2023scaling}. Among the diverse PEFT strategies~\cite{houlsby2019parameter,aghajanyan2020intrinsic,hulora2022,edalati2022krona,wang2022adamix,gheini2021cross,zaken2022bitfit,guo2020parameter,sung2021training,ansell2022composable,lester2021power,li2021prefix,vu2022spot,he2021towards,mao2021unipelt,karimi2021compacter,liu2022few,sung2022lst,chen2022parameter,jia2022visual,chen2022adaptformer,zhang2022neural,jie2023fact,lian2022scaling,luo2023towards}, those based on reparameterization~\cite{aghajanyan2020intrinsic,hulora2022,edalati2022krona} are especially pertinent here. The LoRA technique~\cite{hulora2022} modifies a pretrained weight matrix through the introduction of a product of two low-rank matrices, yielding notable success on NLP benchmarks. While LoRA employs a fixed rank for the inserted matrices, subsequent works~\cite{zhang2022adaptive,valipour2022dylora,zhang2023increlora} have developed approaches to dynamically tune the ranks at different layers, optimizing parameter usage for task requirements. The intuitive nature of low-rank reparameterization has fueled widespread adoption~\cite{dettmers2023qlora,zi2023delta,chavan2023one}. Drawing inspiration from the role of hyperspherical energy in understanding generalization~\cite{liu2018learning,liu2021orthogonal}, the orthogonal finetuning (OFT) method presented by~\cite{qiu2023controlling} offers an alternative by optimizing an orthogonal transformation within each layer, leading to improved generalization and more consistent convergence relative to LoRA in text-to-image diffusion tasks. However, this gain often comes with a higher number of trainable parameters compared to LoRA, highlighting the need for more parameter-efficient variants of OFT. Furthermore, the utility of OFT outside the original application domain remains unexplored. The BOFT approach advances OFT by introducing butterfly factorization, thereby reducing its parameter footprint and allowing demonstration of orthogonal finetuning benefits across broader adaptation scenarios.

\textbf{Butterfly structures}. The radix-2 Cooley-Tukey method breaks down an $N$-point discrete Fourier transform into two smaller $\frac{N}{2}$-point DFTs recursively, resulting in a characteristic butterfly configuration that can be represented as a product of several sparse matrices (collectively referred to as a butterfly matrix). Such butterfly matrices have seen application in parameterizing orthogonal matrices to bypass pivoting in Gaussian elimination and enhance computational efficiency~\cite{parker1995random}, in reinforcing the stability of recurrent neural network training~\cite{jing2017tunable}, and in approximating kernels~\cite{munkhoeva2018quadrature}. Techniques for learning rapid linear transformations with butterfly parameterizations have been developed in~\cite{dao2019learning,dao2019kaleidoscope}. Additionally,~\cite{chen2022pixelated,dao2022monarch} exploit butterfly matrices to facilitate efficient neural network training. Their structure has also proven valuable for fast matrix-vector multiplications~\cite{michielssen1996multilevel,o2010algorithm}, data-sparse matrix representation~\cite{li2015butterfly}, and optimized network communications~\cite{klasing1994broadcasting,li2003linear,rajasingh2016transmission}. Unlike earlier studies, our work leverages butterfly structures specifically to increase the parameter efficiency of OFT.

\section{An Information Transmission View on Orthogonal Finetuning}

We begin by outlining some foundational aspects of OFT. In OFT, the pretrained weight matrix is fine-tuned by expressing the updated matrix as the product of a learnable orthogonal matrix and the original frozen pretrained weight matrix. Unlike LoRA, which modifies the weights by adding a low-rank matrix, OFT relies on a multiplicative update via an orthogonal matrix. LoRA achieves parameter efficiency through its low-rank design; in contrast, the initial formulation of OFT adopts a block-diagonal orthogonal structure~\cite{qiu2023controlling}, where reducing the size of each diagonal block enhances parameter efficiency. Figure~\ref{fig:comp_lora} visually contrasts these approaches. The use of orthogonal transformations in weight fine-tuning is motivated by the goal of preserving the angular relationships between neurons~\cite{liu2018learning,liu2021orthogonal,qiu2023controlling}, thereby maintaining a substantial amount of the semantic information acquired during pretraining.

More specifically, OFT trains an orthogonal matrix $\bm{R}\in\mathbb{R}^{d\times d}$ for a given pretrained linear layer $\bm{W}^0\in\mathbb{R}^{d\times n}$. The standard forward computation, $\bm{z}=(\bm{W}^0)^\top\bm{x}$, is thus modified to $\bm{z}=(\bm{R}\bm{W}^0)^\top\bm{x}$, with $\bm{x}\in\mathbb{R}^{d}$ denoting the input and $\bm{z}\in\mathbb{R}^{n}$ the resulting output. To support zero initialization, $\bm{R}$ is initially set as the identity matrix. The orthogonality of $\bm{R}$ is rigorously maintained during training by leveraging the Cayley parameterization, following~\cite{qiu2023controlling,liu2021orthogonal}: $\bm{R}=(\bm{I}+\bm{Q})(\bm{I}-\bm{Q})^{-1}$, where $\bm{Q}$ is a skew-symmetric matrix satisfying $\bm{Q}=-\bm{Q}^\top$. To promote parameter efficiency, the orthogonal matrix $\bm{R}$ adopts a block-diagonal structure, expressed as $\text{diag}(\bm{R}_1,\bm{R}_2,\cdots,\bm{R}_r)$ where each $\bm{R}_i\in\mathbb{R}^{b\times b}$ is a smaller orthogonal matrix and $br=d$.

While this block-diagonal scheme conserves parameters, it also imposes a structural constraint—namely, that the vector dimensions associated with each neuron (that is, the columns of $\bm{W}^0$) are partitioned into $r$ distinct groups, with each group being independently transformed by its own orthogonal matrix. Although experimental evidence supports the effectiveness of this configuration, there is little rationale for dividing a neuron's dimensions into $r$ groups solely by their indices, pointing to the potential value of dense orthogonal matrices. This consideration leads naturally to a central question: \emph{Is it possible to develop a dense orthogonal matrix that retains parameter efficiency?}

In response to this question, we suggest expressing a dense orthogonal matrix as a sequence of several sparse orthogonal matrices. To unify and clarify the examination of sparsity structures in orthogonal matrix decompositions, we reinterpret the generation of a dense orthogonal matrix within OFT through the lens of information transmission. Specifically, the process of constructing a dense matrix $\bm{R}\in\mathbb{R}^{d\times d}$ as a product $\bm{R}=\bm{B}_m\bm{B}_{m-1}\cdots\bm{B}_1$, where each $\bm{B}_i$ is a square matrix, can be visualized as information being passed through a network with $d\times (m+1)$ nodes (see Figure~\ref{fig:info_trans}). This information transmission analogy stems from the insight that a dense $d$-by-$d$ matrix defines complete connectivity from one set of $d$ nodes to another. For the matrix $\bm{R}$, a zero entry $R_{ij}$ implies the absence of a path for information from the $j$-th to the $i$-th node, while a nonzero $R_{ij}$ signifies possible transmission. As such, breaking down the dense matrix $\bm{R}$ into the sequence $\bm{B}_m\bm{B}_{m-1}\cdots\bm{B}_1$ equates to performing stepwise information exchange in accordance with the graph structures of each $\bm{B}_i$. The information propagates sequentially, first through $\bm{B}_1$ and concluding with $\bm{B}_n$. 

To illustrate, Figure~\ref{fig:info_trans} shows the example of $\bm{R}=\bm{B}_5\bm{B}_4\bm{B}_3\bm{B}_2\bm{B}_1$, alongside visualizations of the sparsity patterns and the resulting graph. The figure represents the matrix multiplication as an unrolled graph, where each matrix $\bm{B}_i$ defines connectivity between nodes at the $i$-th and $(i+1)$-th layers. Concretely, the $(j_1,j_2)$ entry of $\bm{B}_i$ determines whether a directed edge exists from node $j_2$ in level $i$ to node $j_1$ in level $(i+1)$ (with a zero indicating no edge). For $\bm{B}_5\bm{B}_4\bm{B}_3\bm{B}_2\bm{B}_1$ to form a dense matrix, each initial node must be able to reach all nodes at the sixth level. When only considering the product $\bm{R}=\bm{B}_4\bm{B}_3\bm{B}_2\bm{B}_1$, connecting sources at the first level to receivers at the fifth, it becomes apparent that, for instance, node $1$ at the source cannot transmit to node $3$ at the destination—reflected in a zero entry at position $(3,1)$. This correspondence between induced graph structure and the sparsity pattern persists if we further restrict the product to $\bm{R}=\bm{B}_3\bm{B}_2\bm{B}_1$ or $\bm{R}=\bm{B}_2\bm{B}_1$. In general, to realize a dense matrix $\bm{R}\in\mathbb{R}^{d\times d}$ through $m$ factors, the collection $\bm{B}_m,\ldots,\bm{B}_1$ must collectively produce directed paths on a $d\times (m+1)$ grid. These paths, each connecting adjacent layers, must ensure that every node at the initial level can transmit information to all nodes at the $(m+1)$-th level.

Adopting the information transmission perspective offers a useful lens to interpret the sparsity characteristics seen in the factorization matrices utilized within OFT. As depicted in Figure~\ref{fig:blk_diag}, $\bm{R}$ in the canonical OFT configuration exhibits a block-diagonal arrangement. Although such a structure limits the count of learnable parameters, it fails to yield a fully dense $\bm{R}$. The challenge we address is the synthesis of a dense orthogonal matrix from $m$ constituent sparse orthogonal matrices, utilizing the minimal number of necessary trainable parameters. From the standpoint of information transmission, our main requirements for achieving this aim include: \emph{(i)} \textbf{dense connectivity}, ensuring every initial node connects via at least one path to all terminal nodes, and \emph{(ii)} \textbf{minimal redundant edges}, wherein the total edge count remains as low as possible, conditioned on the orthogonality constraint. It is important to emphasize that orthogonality imposes a nuanced constraint on the permissible edges linking subsequent levels. Concretely, for any matrix $\bm{B}_i$ to possess full rank—a prerequisite for orthogonality—there must be $d$ edges establishing a bijective mapping from all nodes at level $i$ to those at level $(i=1)$. Consequently, each pair of consecutive levels requires a minimum of $d$ edges (for instance, four in the scenario illustrated by Figure~\ref{fig:info_trans}). As these $d$ edges are essential for maintaining orthogonality and lack trainable parameters (e.g., in a $d\times d$ orthogonal matrix, these fixed entries are restricted to $\pm 1$), they should be omitted from the effective edge count. The inherent parameter efficiency of orthogonal matrices as opposed to general dense matrices means that orthogonality induces further dependencies among the edges: for example, in a $2\times 2$ orthogonal matrix, a zero at the $(1,1)$ entry enforces a zero at $(2,2)$, so the omission of one edge can necessitate the absence of another. Accordingly, within any permitted edge configuration, orthogonality may induce the creation or elimination of certain edges. Focusing on the distribution of non-zero elements in the resulting orthogonal matrix, the information transmission perspective is especially apt for OFT, as the main concern lies in the location of non-zero trainable elements within $\bm{R}$—the precise values of these elements are inconsequential.

Establishing a straightforward dense linkage between two layers entails $\mathcal{O}(d^2)$ connections (that is, corresponding to a single dense orthogonal matrix), resulting in $d^2-d$ edges—for instance, when $d=4$, there are $12$ such connections. An illustration in Figure~\ref{fig:info_trans} shows a potential matrix decomposition utilizing a total of $10$ edges, which is, in fact, fewer than that required for one dense orthogonal matrix. This perspective provides a graphical approach for examining the parameter efficiency of OFT, making it straightforward to design viable factorizations. Our method is inspired by the notable connectivity scheme from the Cooley-Tukey algorithm, namely butterfly graphs~\cite{cooley1965algorithm}, which can realize all-to-all connectivity between $d$ inputs and $d$ outputs using just $\mathcal{O}(d\log d)$ connections. For context, the arrangement shown in Figure~\ref{fig:info_trans} necessitates $10$ edges to create dense interactions, but the butterfly architecture can accomplish this with merely $8$ edges. In the following, we detail the way in which the butterfly structure enhances parameter efficiency.

\section{Orthogonal Parameterization by Butterfly Factorization}

The butterfly configuration originally arises from the Cooley-Tukey algorithm for efficient computation of the fast Fourier transform. Notably, in a Fourier transform, a minor adjustment in the frequency space can lead to wide-reaching changes in the spatial domain, reflecting a concept similar to our issue of information propagation—every node in the initial layer can transfer information to every node in the final layer. The butterfly structure is also widely adopted in computer networks~\cite{solihin2015fundamentals,li2011linear} for rapid and efficient data exchange. Let $k\geq 2$ be a power of $2$. We first define the \emph{butterfly factor} $\bm{B}^F(k)$ as follows:

\begin{equation}\label{eq:but_factor}
\begin{aligned}
    \bm{B}^F(k)=\begin{bmatrix}
    \text{diag}(\bm{d}_1) & \text{diag}(\bm{d}_2) \\
    \text{diag}(\bm{d}_3) & \text{diag}(\bm{d}_4)
    \end{bmatrix}\in\mathbb{R}^{k\times k},
\end{aligned}
\end{equation}

where each $\bm{d}_i\in\mathbb{R}^{\frac{k}{2}}$ for all $i$ are vectors of length $\frac{k}{2}$. For $d=2^N$, the $d$-dimensional \emph{butterfly matrix} $\bm{B}(d)\in\mathbb{R}^{d\times d}$ is then recursively specified by

\begin{equation}
\begin{aligned}
    \!\!\bm{B}(d)=\tilde{\bm{B}}(d,d)\!\cdot\!\begin{bmatrix}
    \bm{B}_1(\frac{d}{2})\!\!\!\!\! & \bm{0} \\
    \bm{0}\!\!\!\!\! &\bm{B}_2(\frac{d}{2})
    \end{bmatrix}= \tilde{\bm{B}}(d,d)\tilde{\bm{B}}(d,\frac{d}{2})\cdots\tilde{\bm{B}}(d,2),
\end{aligned}
\end{equation}

where $\bm{B}_1(\frac{d}{2})$ and $\bm{B}_2(\frac{d}{2})$ represent two butterfly matrices of dimension $\frac{d}{2}$. We introduce the \emph{butterfly component} as $\thickmuskip=2mu \medmuskip=2mu \tilde{\bm{B}}(d,k)=\text{diag}(\bm{B}^F_1(k),\cdots,\bm{B}^F_{d/k}(k))$, which constructs a block-diagonal matrix of size $d\times d$ composed of blocks each of size $k$. Here, each $\bm{B}^F_i(k)$ for every $i$ denotes a butterfly factor as specified in Equation~\ref{eq:but_factor}. With this setup, we are prepared to employ butterfly matrices as a means of parameterizing orthogonal matrices. This is feasible provided that all butterfly matrix factors $\tilde{\bm{B}}(d,k)$, for all relevant $k$, are orthogonal. We begin by analyzing the block-diagonal matrix $\tilde{\bm{B}}(d,2)$ where the block size is $2$. Orthogonality for $\tilde{\bm{B}}(d,2)$ can be straightforwardly enforced using the Cayley transform or by representing each $2$-dimensional block as a rotation, consistent with approaches from \cite{liu2021orthogonal,qiu2023controlling}. Each butterfly component’s non-zero structure can be considered as a permutation of the non-zero structure in $\tilde{\bm{B}}(d,2)$, enabling all butterfly components to be systematically parameterized as orthogonal matrices. Consequently, this structure offers an efficient approach to parameterizing orthogonal matrices via compositions of many $2\times 2$ orthogonal matrices. Extending the butterfly construction according to \cite{chen2022pixelated}, we define a block butterfly component $\tilde{\bm{B}}^b(d,k)$, in which each entry of $\bm{d}_i$ is generalized to be a $b\times b$ matrix. To enforce orthogonality for $\tilde{\bm{B}}^b(d,2)$, each $2b\times 2b$ block is parameterized as an orthogonal matrix. The remaining butterfly components $\tilde{\bm{B}}^b(d,k)$, with $k>2$, maintain a non-zero pattern that is a blockwise permutation of the pattern found in $\tilde{\bm{B}}^b(d,2)$, thus they too can be parameterized as orthogonal matrices in the same way. Combining these concepts, the forward computation in BOFT can be written as

\begin{equation}
    \begin{aligned}\nonumber
    \bm{z}=\big(\bm{R}(m,b)\cdot\bm{W}^0\big)^\top\bm{x},~~~~\text{s.t.}~\bigg{\{}\bm{R}(m,b)=\prod_{i=1}^m \tilde{\bm{B}}^b_{(d,i)}~~\&~~\underbrace{\big(\tilde{\bm{B}}^b_{(d,j)}\big)^\top\tilde{\bm{B}}^b_{(d,j)}=\tilde{\bm{B}}^b_{(d,j)}\big(\tilde{\bm{B}}^b_{(d,j)}\big)^\top\!=\bm{I}_d}_{\forall j\in [1, m]}\bigg{\}},
    \end{aligned}
\end{equation}

Here, we use $\tilde{\bm{B}}^b(d,2^{m - i+1})$ and abbreviate it as $\tilde{\bm{B}}^b_{(d,i)}$, with $\bm{I}_d$ denoting the $d \times d$ identity matrix. The orthogonal matrix $\bm{R}(m,b)\in\mathbb{R}^{d\times d}$ results from sequentially multiplying several orthogonal butterfly blocks. For notational simplicity, we define BOFT utilizing $\bm{R}(m,\frac{b}{2})$ as BOFT($m$,$b$) for $b\geq2$. Special cases arise for certain values of $m$ and $b$: if $\thickmuskip=2mu \medmuskip=2mu m=1$, then BOFT($1$,$b$) collapses to the block-diagonal OFT~\cite{qiu2023controlling} where each block is of size $b$; in the case $b=d$, BOFT($1$,$d$) matches the original OFT~\cite{qiu2023controlling} with a dense unconstrained orthogonal transformation. When $m = \log \frac{2d}{b}$, choosing the block butterfly matrix $\bm{B}^{\frac{b}{2}}(d)$ as $\bm{R}$ recovers a dense orthogonal matrix $\bm{R}$. Generally, BOFT($m$,$b$) introduces $\frac{1}{2}(b-1)dm$ effective learnable parameters when fine-tuning a linear transformation of size $d\times n$. For the specific scenario employing the butterfly structure (that is, $m=\log d, b=2$), the parameter cost is $\mathcal{O}(d\log d)$. Conversely, the standard OFT with a dense orthogonal matrix has complexity $\mathcal{O}(d^2)$, and the block-diagonal variant with $r$ blocks requires $\mathcal{O}(bd)$ parameters. Therefore, to obtain a dense orthogonal matrix in the conventional OFT, the block size must be set to $b=d$, whereas BOFT allows for flexibility in the choice of $b$ for the same effect.

\textbf{Identity initialization for BOFT}. A common approach in fine-tuning is to initialize the finetuned parameters to match the pretrained weights as closely as possible. For instance, LoRA typically applies zero initialization to its low-rank updates. In BOFT, we initialize every butterfly factor as an identity matrix (specifically, by initializing the underlying skew-symmetric matrices with zeros when employing the Cayley transform).

\textbf{Multiplicative Dropout for BOFT}. The Dropout mechanism in LoRA~\cite{hulora2022}, inspired by the classic Dropout~\cite{srivastava2014dropout}, helps guard against overfitting in the low-rank adaptation. Standard Dropout is directly compatible with LoRA’s additive formulation, but it is not suitable for BOFT because BOFT relies on multiplicative modifications. To mitigate this, we introduce a multiplicative Dropout specifically tailored for BOFT: at each application, a random selection of $p_1\%$ of the butterfly layers and $p_2\%$ of the blocks within those layers are chosen and replaced with identity blocks, leveraging the fact that $\bm{R}(m,b)$ is a product of $m$ butterfly factors, each of which can be permuted into $2b \times 2b$ block-diagonal orthogonal submatrices.

\section{Intriguing Insights and Discussions}

\textbf{Expressivity of BOFT}. While the butterfly structure combined with permutations is capable of exactly representing various well-known fast linear transforms~\cite{dao2019learning,dao2019kaleidoscope} (\emph{e.g.}, fast Fourier transform, Hadamard transform), the extent to which an orthogonal butterfly matrix can approximate an arbitrary orthogonal matrix is still an open question. To examine this, we perform an empirical study attempting to approximate a random dense orthogonal matrix~\cite{anderson1987generation} of dimension $512\times 512$. In Figure~\ref{fig:para_eff}, the reported values are averaged across 10 random initializations. The y-axis corresponds to the approximation error, while the x-axis reflects the number of effective trainable parameters. Each curve with a particular color represents BOFT with a fixed block size, and the leftmost point on a curve shows the error for BOFT($1$,$b$) (\emph{i.e.}, standard block-diagonal OFT with block size $b$). Across the board, BOFT demonstrates greater parameter efficiency compared to OFT. As an illustration, BOFT($9$,$2$) achieves superior expressivity relative to BOFT($1$,$16$) but with significantly fewer parameters. Typically, using a smaller $b$ and a larger $m$ yields more efficient parameter usage; for instance, BOFT($6$,$4$) attains similar approximation errors as BOFT($2$,$16$) while utilizing far fewer parameters. Broadly, the butterfly matrix imposes more structural constraints within the orthogonal group than the block-diagonal form, which makes BOFT verifiably more expressive than OFT when matched on block size.

\begin{theorem}[Expressivity of BOFT]\label{thm:exp_boft}
    BOFT is strictly more expressive than OFT at identical block sizes. In order to cover all orthogonal matrices of size $d$, it suffices to compose butterfly matrices as $\bm{B}_{d-1,1}(d)\bm{B}^\top_{d-1,2}(d)\cdots\bm{B}_{1,1}(d)\bm{B}^\top_{1,2}(d)$, where each $\bm{B}_{i,j}(d)$ is a butterfly matrix.
\end{theorem}

The result in Theorem~\ref{thm:exp_boft} points to a direct path for generalizing BOFT. Specifically, we can consider an extended formulation: $\thickmuskip=2mu \medmuskip=2mu \bm{R}^G(m_1,b_1,m_2,b_2,l)=\bm{R}_{l,1}(m_1,b_1)\bm{R}_{l,2}^T(m_2,b_2)\cdots\bm{R}_{1,1}(m_1,b_1)\bm{R}_{1,2}^T(m_2,b_2)$, where each $\bm{R}_{i,j}^T(m,b)$ denotes the orthogonal component in BOFT. By setting $m_1=m_2=\log d$, $b_1=b_2=2$, and $l=d-1$, $\bm{R}^G(m_1,b_1,m_2,b_2,l)$ spans the entire orthogonal group—this construction aligns with what is called the kaleidoscope hierarchy~\cite{dao2019kaleidoscope}. However, it is important to recognize that increased expressivity does not necessarily yield better results during finetuning. In particular, despite the capacity of full finetuning for universal representation, it often performs poorly in practice. Thus, balancing expressivity against regularity is crucial for ensuring good generalization in model finetuning. The BOFT framework enhances the finetuning parameter space with embedded structural priors, facilitating the identification of an optimal balance between model expressiveness and regularization.

\textbf{Spectral properties}. Compared to LoRA, orthogonal finetuning typically achieves superior spectral characteristics, principally due to its exact preservation of the spectral norm of the original pretrained weight matrix $\bm{W}^0$. This property is evident via the singular value decomposition: expressing $\bm{W}^0$ as $\bm{U}\bm{\Sigma}\bm{V}^\top$—with $\bm{U}, \bm{V}$ as orthogonal matrices and $\bm{\Sigma}$ as a diagonal matrix containing the singular values. Both OFT and BOFT effectuate finetuning by left-multiplying an orthogonal matrix $\bm{R}$, resulting in updated weights $\bm{R}\bm{U}\bm{\Sigma}\bm{V}^\top$, a transformation that leaves the largest singular value (i.e., the spectral norm of $\bm{W}^0$) unchanged. Maintaining this spectral norm has been demonstrated to significantly enhance training robustness and generalizability~\cite{miyato2018spectral,yoshida2017spectral}. Additional theoretical properties are provided in Appendix~\ref{app:sn_property}.

\textbf{Orthogonal finetuning as learning bilinear similarity}. The BOFT method can be reformulated as optimizing for a bilinear similarity function, $\bm{w}^0_i\bm{R}\bm{x}$, where $\bm{w}^0_i$ refers to the $i$-th neuron (i.e., a column vector) from $\bm{W}^0$. In this context, BOFT seeks to learn the bilinear similarity matrix $\bm{R}$, subject to the orthogonality constraint on $\bm{R}$, thus establishing a close link with both distance metric learning~\cite{xing2002distance} and the theory of bilinear forms~\cite{roman2005advanced}.

\textbf{Inductive bias and generalization in BOFT}. The matrix $\bm{R}(m,b)$ in BOFT is typically constrained to a structured subset within the group of orthogonal matrices, effectively reducing the expressiveness of the hypothesis space. This structure inherently induces an inductive bias in the learning process. We posit that the inductive bias imposed by the butterfly factorization structure in BOFT benefits model generalization, as it emulates the structural motifs present in a number of classic linear transforms~\cite{dao2019kaleidoscope}, such as the discrete Fourier, sine/cosine, and Hadamard transforms. Additionally, the sparse nature of the matrix factorization adopted by BOFT is likely to introduce further implicit inductive biases~\cite{gunasekar2017implicit,li2018algorithmic,liu2019neural}.

\textbf{Comparison to butterfly-based sparse training}. Several previous studies~\cite{dao2019learning,dao2019kaleidoscope,chen2022pixelated,dao2022monarch} have investigated sparse training methods leveraging butterfly parameterization. These approaches predominantly emphasize expressing weight matrices directly using butterfly parameterization, training neural networks from the initial state. In contrast, \cite{dao2022monarch} explores fine-tuning by first projecting pretrained weights onto a specialized form of butterfly matrices and subsequently optimizing only those projected components for downstream objectives. In BOFT, we introduce an alternative fine-tuning methodology, which modifies weights utilizing weight matrices that are shared across layers.

\input{rebuttal-results}

\section{Concluding Remarks and Limitations}

In this work, we introduce Orthogonal Butterfly, a general and parameter-efficient fine-tuning framework for foundation models that is built on the butterfly matrix structure. The principal idea behind improved parameter efficiency is to represent a dense orthogonal matrix as the product of several sparse orthogonal matrices. To facilitate the identification of valid matrix factorizations, we present a graphical approach grounded in information transmission principles. Within this paradigm, the butterfly architecture emerges as a practical means to realize sparse orthogonal matrix factorization. We provide empirical evidence that BOFT achieves strong performance when applied to fine-tuning across large-scale language models, vision models, and text-to-image generation architectures. Our results also corroborate the overall effectiveness of BOFT as a flexible fine-tuning strategy.

Nevertheless, BOFT exhibits some limitations. Because its final orthogonal matrix results from composing multiple orthogonal matrices, training with BOFT incurs somewhat greater computational overhead compared to OFT. Addressing this increased runtime during training remains an open challenge. Nonetheless, once fine-tuning concludes, the orthogonal matrices produced by BOFT can be incorporated into the pretrained model directly, ensuring that inference speed is unaffected. Additionally, it is not yet clear whether the butterfly topology constitutes the most effective mechanism for information transfer. Our information transmission-based framework further suggests analogies with the field of computer networking, where the problem of efficiently structuring network topologies for data transmission is well-explored. We anticipate that pursuing more efficient network topologies could further advance the composition of dense orthogonal matrices.